\section{Discussion}
\label{sec_discussion}

\acresetall
% The \ac{mem} shape shown in \cref{fig_no_circle} has a curvature which varies along the \ac{mem} profile, and thus differs from the classical circle-arc obtained at the interface between two fluids \cite{landauFluidMechanics1987,degennesCapillarityWettingPhenomena2004}.

We developed a numerical study on the interaction between a \ac{bflu} and a \ac{mem}. Our analysis is based on the \ac{fe} method, and it leverages the \ac{fenics} \cite{loggAutomatedSolutionDifferential2012}, combined with \ac{irene} \cite{worthmullerIRENEFluIdLayeR2025}.
In order to guide and introduce the reader to the interaction problem between a \acl{bflu} and \acl{mem}, we provided two introductory examples of fluid-structure interaction \cite{kamenskyLectureNotesMAE}, in which the \acl{bflu} interacts with a \acl{rig} and with an \acl{ela}, see \cref{sec_fl_rigid,sec_fl_elastic}. These Sections allow to introduce the necessary physical an mathematical quantities, as well as to prove the  stability of our numerical scheme.

%sign